% Generated by experimento/render_final_tables.py from audited R2/A1 data.
\begin{table}[htbp]
\centering
\caption{Diferenças pareadas do p95 HTTP: assíncrona menos síncrona, em ms}
\label{tab:final-pares}
\small
\setlength{\tabcolsep}{4pt}
\begin{tabular}{lrrrrr}
\hline
Condição & Bloco 1 & Bloco 2 & Bloco 3 & Bloco 4 & Bloco 5 \\ \hline
Constante 12/s & -2,26 & -2,80 & -1,77 & -4,54 & -2,28 \\
Constante 24/s & -1,59 & -1,76 & -2,44 & -1,60 & -3,20 \\
Constante 36/s & -1,53 & -1,67 & -1,60 & -1,59 & -3,47 \\
Constante 48/s & -1,90 & -1,60 & -1,98 & -1,44 & -1,75 \\
Rajada & -0,24 & -1,60 & -1,72 & -1,25 & -2,56 \\
Banco parado & -17,01 & -32,12 & -2,90 & -20,08 & -4,54 \\
\hline
\end{tabular}
\fonte{Elaborado pelo autor a partir da campanha de 27/10/2026.}
\vspace{0.2cm}
\begin{minipage}{\textwidth}
\footnotesize
Cada coluna é um par com o mesmo conjunto de dados, não uma amostra de evento. Diferenças negativas indicam menor p95 HTTP da assíncrona nestes cinco pares. A análise é descritiva; as diferenças e dispersões completas estão no conjunto identificado no Apêndice~\ref{ap:rastreabilidade}. Não existe par síncrono para consumidor parado.
\end{minipage}
\end{table}
